\chapter{不想教的東西}

本章討論應該避免的 C 語言特性和實踐。 瞭解這些反面教材，有助於寫出更好、更安全的程式。

\section{goto 的問題}

\texttt{goto} 無條件跳轉到標籤。 雖然在某些特殊情況下有用，但容易導致 spaghetti code。

\subsection*{goto 導致的混亂}

\begin{code}{c}{goto spaghetti code}
// WRONG: hard to follow
int process(int n) {
    int result = 0;
    if (n < 0) goto error;
    if (n == 0) goto done;

    result = n * 2;
    if (result > 100) goto large;

    printf("Small: %d\n", result);
    goto done;

large:
    printf("Large: %d\n", result);
    goto done;

error:
    printf("Error\n");

done:
    return result;
}
\end{code}

\subsection*{改用結構化控制流}

\begin{code}{c}{structured control flow}
// BETTER: clear flow
int process(int n) {
    if (n < 0) {
        printf("Error\n");
        return -1;
    }

    if (n == 0) {
        return 0;
    }

    int result = n * 2;

    if (result > 100) {
        printf("Large: %d\n", result);
    } else {
        printf("Small: %d\n", result);
    }

    return result;
}
\end{code}

\subsection*{goto 的稀有用途}

在某些情況下 \texttt{goto} 能簡化代碼，如錯誤處理清理。

\begin{code}{c}{goto for cleanup}
#include <stdio.h>
#include <stdlib.h>

int process_file(const char *filename) {
    FILE *f = NULL;
    int *data = NULL;
    int result = 0;

    f = fopen(filename, "r");
    if (!f) goto error;

    data = (int*)malloc(1000 * sizeof(int));
    if (!data) goto error;

    // Process file...
    result = 1;

    goto cleanup;

error:
    result = -1;

cleanup:
    if (data) free(data);
    if (f) fclose(f);
    return result;
}
\end{code}

\keypoint{\texttt{goto} 通常應避免。 只在錯誤清理時考慮使用。}

\section{全域變數的陷阱}

全域變數導致隱藏的依賴性、難以除錯、併發問題。

\subsection*{隱藏的依賴}

\begin{code}{c}{global state problem}
// WRONG: implicit dependency
int counter = 0;

void increment() {
    counter++;
}

int get_value() {
    return counter * 2;
}

int main() {
    increment();
    printf("%d\n", get_value());  // depends on counter
    // Hard to trace where counter is modified
    return 0;
}
\end{code}

\subsection*{參數傳遞的清晰性}

\begin{code}{c}{explicit parameters}
// BETTER: clear dependencies
int increment(int counter) {
    return counter + 1;
}

int get_value(int counter) {
    return counter * 2;
}

int main() {
    int counter = 0;
    counter = increment(counter);
    printf("%d\n", get_value(counter));
    return 0;
}
\end{code}

\subsection*{何時用 static 而非全域}

\begin{code}{c}{static for encapsulation}
// Module-level counter (file scope)
static int local_counter = 0;

void increment() {
    local_counter++;
}

int get_counter() {
    return local_counter;
}

// Outside this file cannot access local_counter directly
\end{code}

\subsection*{併發問題}

\begin{code}{c}{race condition}
// DANGER: race condition in multi-threaded code
int shared_count = 0;

void worker_a() {
    for (int i = 0; i < 1000; i++) {
        shared_count++;  // not atomic
    }
}

void worker_b() {
    for (int i = 0; i < 1000; i++) {
        shared_count++;  // race condition
    }
}

// Result: shared_count might be < 2000
// Use mutex or atomic operations in multi-threaded code
\end{code}

\keypoint{全域變數導致難以追蹤的依賴。 優先用參數傳遞。}

\section{過度最佳化}

過早最佳化是萬惡之源。 正確性和可讀性優先。

\subsection*{不必要的複雜性}

\begin{code}{c}{over optimization}
// WRONG: overcomplicated
int fibonacci_bad(int n) {
    return n <= 1 ? n :
           ((n & 1) ? 
            ((1 << (n-1)) * fibonacci_bad(n-1) >> 1) :
            fibonacci_bad(n >> 1) * (2*fibonacci_bad(n>>1)+1));
    // Hard to read, hard to maintain, not necessarily faster
}

// BETTER: clear and correct
int fibonacci(int n) {
    if (n <= 1) return n;
    int a = 0, b = 1;
    for (int i = 2; i <= n; i++) {
        int temp = a + b;
        a = b;
        b = temp;
    }
    return b;
}
\end{code}

\subsection*{正確的最佳化流程}

\begin{code}{text}{optimization process}
1. Write clear, correct code first
2. Test thoroughly
3. Measure performance (profiling)
4. Identify bottlenecks
5. Optimize only the hot spots
6. Test again to ensure correctness
7. Measure improvement
\end{code}

\subsection*{效能測量}

\begin{code}{c}{performance measurement}
#include <stdio.h>
#include <time.h>

void measure_function(int iterations) {
    clock_t start = clock();

    for (int i = 0; i < iterations; i++) {
        fibonacci(30);  // or other function
    }

    clock_t end = clock();
    double elapsed = (double)(end - start) / CLOCKS_PER_SEC;
    printf("Time: %.3f seconds\n", elapsed);
}

int main() {
    measure_function(100);
    return 0;
}
\end{code}

\subsection*{權衡點}

\begin{code}{c}{performance tradeoffs}
// Fast but hard to read
int optimized_version(int x) {
    return (x * 3) >> 1;  // multiply by 1.5
}

// Slower but clear
int readable_version(int x) {
    return (x * 3) / 2;
}

// If this function is called rarely, use readable version
// Only optimize if profiler shows it's a bottleneck
\end{code}

\keypoint{最佳化時應基於測量。 可讀性優先，除非效能是關鍵。}

\section{總結：好的實踐}

\subsection*{應該做的事}

\begin{itemize}
\item 寫清晰、可讀的代碼
\item 使用函式和參數傳遞而不是全域變數
\item 檢查所有返回值和指標
\item 遵循命名約定
\item 編寫單元測試
\item 使用版本控制
\item 為複雜邏輯編寫文件
\end{itemize}

\subsection*{應該避免的事}

\begin{itemize}
\item 不要用 \texttt{goto}（除非必要）
\item 不要過度使用全域變數
\item 不要過早最佳化
\item 不要忽視記憶體管理
\item 不要假設編譯器會修復邏輯錯誤
\item 不要在不測量的情況下宣稱效能改進
\end{itemize}

\keypoint{好的程式是平衡的。 正確性、可讀性、效能都重要，但優先順序應合理。}
